% !TEX root = ../aa_SoM_Chapters.tex

% ö = \%c3\%b6
% ä = \%C3\%A4

\xdef\sectiontitle{\Isectiontitle} %aiheuttama

%%%%%%%%%%%%%%%
\begin{frame}[label=1_8_a]%[label=1_8_TOC1]
\frametitle{\texorpdfstring{\culine[\sectiontitleulinecolor]{\sectiontitle}}{\sectiontitle} }

\vspace{0.5cm}

\begin{itemize}[leftmargin=0.5cm,itemsep=2mm]
    \item[1.1]<1-> \hyperlink{1_1_a}{\IaSubsectiontitle}
    \item[1.2]<1-> \hyperlink{1_2_a}{\IbSubsectiontitle}	
    \item[1.3]<1-> \hyperlink{1_3_a}{\IcSubsectiontitle}	
    \item[1.4]<1-> \hyperlink{1_4_a}{\IdSubsectiontitle}
    \item[1.5]<1-> \hyperlink{1_5_a}{\IeSubsectiontitle}
    \item[1.6]<1-> \hyperlink{1_6_a}{\IfSubsectiontitle}
    \item[1.7]<1-> \hyperlink{1_7_a}{\IgSubsectiontitle}
\end{itemize}

\vspace{-1cm}
\begin{itemize}[leftmargin=9.5cm,itemsep=2mm]
    \item[1.8]<1-> \hyperlink{1_8_a}{\CDAlert[red]{\IhSubsectiontitle}}
	\item[1.9]<1-> \hyperlink{1_9_a}{\IiSubsectiontitle}
	\item[1.10]<1-> \hyperlink{1_10_a}{\IjSubsectiontitle}
	\item[1.11]<1-> \hyperlink{1_11_a}{\IkSubsectiontitle}
\end{itemize}

%\endofslide 
\end{frame}
%%%%%%%%%%%%%%%


\xdef\sectiontitle{\IhSubsectiontitle} 

%%%%%%%%%%%%%%%
\begin{frame}
\frametitle{\texorpdfstring{\culine[\sectiontitleulinecolor]{\sectiontitle}}{\sectiontitle} }

\begin{itemize}
	\item Unlike electrons (and He-4 atoms), photons are bosons with spin $=1$ 
	%\item The behavior of  
	\item To describe photon gas\appear{, we need to know their \CDAlert[fuchsia]{density of states}}

\item The energy of photons is linked to frequency $f$ \appear{(instead of a quantum number $n$ of energy level):} $ \qquad \appear{E =h f } \equal \appear{h \frac{c}{\lambda}} $
\end{itemize}

\newcommand\myeq{\stackrel{\mathclap{\footnotesize\mbox{analogous to}}}{~\Leftrightarrow~}}
\begin{minipage}{8cm}
\begin{itemize}

	\item First, consider 3D-quantum well of width $L$\appear{, where quanta with $\lambda=2 L$ (at minimum) can fit}
	\[
\begin{aligned}
n_{x}^{2} \plus \appear{n_{y}^{2}} \plus \appear{n_{z}^{2}}  \equal  \appear{n^{2}}  \qquad \myeq \qquad  \appear{R^{2}}  \equal  \appear{ \bigg( \Frac{L}{\lambda / 2}  \bigg)^{2}}
\end{aligned}
\]
\item[] allowing us to calculate the radius in phase space:
\[
R  \equal  \frac{L}{\lambda / 2}  \equal  \frac{2 L}{\lambda}  \equal  \frac{2 L E}{h c}
\]
\end{itemize}
\end{minipage}


\xdef\loc{south east}
\begin{tikzpicture}[remember picture,overlay] % anchor=south
    \node (A) [yshift=-0.25*\figYshift,xshift=\figXshift, anchor=\loc] at (current page.\loc) {\includegraphics[page=1,width=5.2cm]{Figures_booklet/SOM_Tipler_Statistics_and_Bonding_figures/SOM_Tipler_figure.png}}; %scale=\figscale. % 8cm max hei
\end{tikzpicture}


\endofslide 
%\endofsection
\end{frame}
%%%%%%%%%%%%%%%

%%%%%%%%%%%%%%%
\begin{frame}
\frametitle{\texorpdfstring{\culine[\sectiontitleulinecolor]{\sectiontitle}}{\sectiontitle} }

\begin{itemize}
	\item How many 'lattice points' \appear{can we have in a $n$-space volume with radius $R$?}
\[
\begin{aligned}
& & R& \equal \frac{2 L E}{h c} \\ 
&\appear{\Rightarrow} \qquad& \appear{N} &\equal \appear{2 \cdot} \frac{1}{8} \appear{\cdot} \frac{4}{3} \appear{\pi R^{3}}  \equal  \frac{\pi}{3} \appear{\textrm{((} \frac{2 L E}{h c} \textrm{))} ^{3}}  
\equal  \frac{8 \pi}{3} \frac{E^{3}}{h^{3} c^{3}} \appear{L^{3}}  
\equal \frac{8 \pi}{3} \frac{E^{3}}{h^{3} c^{3}} \appear{V}
\end{aligned}
\]

\item[] where $N$ is the number of states within a sphere of radius $R=2 L E /(h c)$ 

\item Now we can calculate the number of points within an infinitesimal energy range $d E$:
%With the obtained number of points $N$ \appear{within a sphere of radius $R=2 L E /(h c)$ in the $n$-space}\appear{, we can derive how many points we have per infinitesimal energy range $d E$:}
\[
\frac{d N}{d E}  \equal \frac{8 \pi}{3} \frac{3 E^{2}}{h^{3} c^{3}} \appear{V}
 \equal \frac{8 \pi E^{2}}{h^{3} c^{3}} \appear{V} 
 \qquad \Rightarrow \qquad 
 \frac{d N}{d f}  \equal \frac{8 \pi f^{2}}{c^{3}} \appear{V} \qquad \appear{(E=h f} \quad \Rightarrow \quad \appear{d E=h d f)}
\]

\end{itemize}

\endofslide 
%\endofsection
\end{frame}
%%%%%%%%%%%%%%%

%%%%%%%%%%%%%%%
\begin{frame}
\frametitle{\texorpdfstring{\culine[\sectiontitleulinecolor]{\sectiontitle}}{\sectiontitle} }

\begin{itemize}[itemsep=2.5mm]
	\item So, to interpret the result\appear{, we have obtained the number of states within a certain frequency range:}
\[
\frac{\text { different number of states in range df }}{d f}  \equal \frac{8 \pi V}{c^{3}} \appear{f^{2}}
\]
\item Or, if energy is considered:
\[
\frac{\text { different number of states in range dE }}{d E / h}  \equal \frac{8 \pi V}{c^{3}} \appear{(E / h)^{2}}
\]

\item Thus, for the \CDAlert[red]{density of states $D=D(E)$} \appear{in a system of photons (in a volume $V$)}\appear{, we obtain:}
\[
D(E)  \equal \frac{8 \pi V}{h^{3} c^{3}} \appear{E^{2}}
\]
\end{itemize}

%\xdef\loc{east}
%\begin{tikzpicture}[remember picture,overlay] % anchor=south
%    \node (A) [yshift=0cm,xshift=\figXshift, anchor=\loc] at (current page.\loc) {\includegraphics[page=1,height=7cm]{Figures_booklet/SOM_9_Statistical_mechanics_figures/SOM_Figure_17.png}}; %scale=\figscale. % 8cm max hei
%\end{tikzpicture}


\endofslide 
%\endofsection
\end{frame}
%%%%%%%%%%%%%%%

%%%%%%%%%%%%%%%
\begin{frame}
\frametitle{\texorpdfstring{\culine[\sectiontitleulinecolor]{\sectiontitle}}{\sectiontitle} }

\begin{itemize}[itemsep=1.7mm]
	\item Note, that the photon gas differs from the harmonic oscillator\appear{, or e.\,g. liquid helium case}

\item In the photon gas, we \Alert[blue]{do not actually fix the number of photons beforehand}\appear{, since all systems absorb and emit photons continuously}\appear{, and therefore the temporary number of photons is rather dictated by the physics of the system} \appear{and its conditions}

\item The \CDAlert[red]{Bose–Einstein statistics} for bosons were given by: 
\[ \mathbb{N} (E)_{B E}  \equal  \frac{1}{B e^{E / k_{B} T}-1}
\] 
\item[] where the parameter $B$ should be linked to number of particles $N$ 

\item Since $N$ is not fixed\appear{, a solution is to set $B\equiv\appear{1}$} \appear{(for photon gas)}
\[
 \mathbb{N} (E)_{B E, \text { photon } g a s}=\frac{1}{e^{E / k_{B} T}-1}
\]
\end{itemize}

%\xdef\loc{east}
%\begin{tikzpicture}[remember picture,overlay] % anchor=south
%    \node (A) [yshift=0cm,xshift=\figXshift, anchor=\loc] at (current page.\loc) {\includegraphics[page=1,height=7cm]{Figures_booklet/SOM_9_Statistical_mechanics_figures/SOM_Figure_17.png}}; %scale=\figscale. % 8cm max hei
%\end{tikzpicture}

\endofslide 
%\endofsection
\end{frame}
%%%%%%%%%%%%%%%

%%%%%%%%%%%%%%%
\begin{frame}
\frametitle{\texorpdfstring{\culine[\sectiontitleulinecolor]{\sectiontitle}}{\sectiontitle} }

\begin{itemize}[itemsep=3mm]
	\item \textbf{Q}: Compare quantitatively the internal energies in electromagnetic radiation and matter
\item \textbf{A}: We could calculate the average energy as given before: 
\[
<E>\,  \equal  \frac{\nint E  \mathbb{N} (E) D(E) d E}{\nint  \mathbb{N} (E) D(E) d E}
\]
\item[] where the nominator stands for the energy of the system \appear{(and denumerator stands for total number of particles)} 

\item But, we could as well use only the nominator to obtain the internal energy of EM-radiation\appear{, or 'photon gas'} 

%\item Using $E=k_{B} T x$, we get:
%\[
%\begin{aligned}
%&U_{p h o t o n}=\nint_{0}^{\infty} E  \mathbb{N} (E) D(E) d E=\nint_{0}^{\infty} E \frac{1}{e^{E / k_{B} T}-1} \frac{8 \pi V}{h^{3} c^{3}} E^{2} d E \\
%&=\frac{8 \pi V}{h^{3} c^{3}} \nint_{0}^{\infty} k_{B} T x \frac{1}{e^{x}-1} \textrm{((} k_{B} T x \textrm{))} ^{2} k_{B} T d x=\frac{8 \pi V k_{B}^{4} T^{4}}{h^{3} c^{3}} \nint_{0}^{\infty} \frac{x^{3} d x}{e^{x}-1}=\frac{8 \pi V k_{B}^{4} T^{4}}{h^{3} c^{3}} \cdot \frac{\pi^{4}}{15}
%\end{aligned}
%\]
\end{itemize}

%\xdef\loc{east}
%\begin{tikzpicture}[remember picture,overlay] % anchor=south
%    \node (A) [yshift=0cm,xshift=\figXshift, anchor=\loc] at (current page.\loc) {\includegraphics[page=1,height=7cm]{Figures_booklet/SOM_9_Statistical_mechanics_figures/SOM_Figure_17.png}}; %scale=\figscale. % 8cm max hei
%\end{tikzpicture}


\endofslide 
%\endofsection
\end{frame}
%%%%%%%%%%%%%%%

%%%%%%%%%%%%%%%
\begin{frame}
\frametitle{\texorpdfstring{\culine[\sectiontitleulinecolor]{\sectiontitle}}{\sectiontitle} }

\begin{itemize}
%\item But, we could as well use only the nominator to obtain the internal energy of EM-radiation\appear{, or 'photon gas'} 

\item Using the change of variables, $\appear{E}  \equal \appear{k_{B} T x} \quad \Rightarrow \quad \appear{x}  \equal \appear{E/(k_{B} T)}$, the nominator turns into:
\[
\begin{aligned}
U_{photon}  &\equal  \appear{\nint_{0}^{\infty}} \appear{E  \mathbb{N} (E)} \appear{D(E)} \appear{d E}
 \equal  
 \appear{\nint_{0}^{\infty} E} \frac{1}{e^{E / k_{B} T}-1} \frac{8 \pi V}{h^{3} c^{3}} \appear{E^{2}} \appear{d E} \\
& \equal 
\frac{8 \pi V}{h^{3} c^{3}} \appear{\nint_{0}^{\infty}} \appear{k_{B} T x} \frac{1}{e^{x}-1} \appear{\textrm{((} k_{B} T x \textrm{))} ^{2}} \appear{k_{B} T} \appear{d x}
 \equal 
 \frac{8 \pi V k_{B}^{4} T^{4}}{h^{3} c^{3}} \appear{\nint_{0}^{\infty}} \frac{x^{3} d x}{e^{x}-1}  \\
 &\equal 
 \frac{8 \pi V k_{B}^{4} T^{4}}{h^{3} c^{3}} \appear{\cdot} \frac{\pi^{4}}{15}
\end{aligned}
\]
\end{itemize}


\endofslide 
%\endofsection
\end{frame}
%%%%%%%%%%%%%%%

%%%%%%%%%%%%%%%
\begin{frame}
\frametitle{\texorpdfstring{\culine[\sectiontitleulinecolor]{\sectiontitle}}{\sectiontitle} }

\begin{itemize}
	\item So, for photon gas we obtain 
	\[
	U_{p h o t o n}=\frac{8 \pi V k_{B}^{4} T^{4}}{h^{3} c^{3}} \appear{\cdot} \frac{\pi^{4}}{15}  \equal \frac{8 \pi^{5} V k_{B}^{4} T^{4}}{15 h^{3} c^{3}}
	\]
\item[] At temperature of $300 \mathrm{~K}$ \appear{we obtain a value of:} $\quad \frac{U_{p h o t o n}}{V} \appear{=6.1 \times 10^{-6} \mathrm{~J} / \mathrm{m}^{3}}$

\item For ordinary \appear{(\CDAlert[fuchsia]{monatomic} ideal)} \appear{gas the internal energy at NTP is given by}
\[
\frac{U_{\text {material }}}{V}
 \equal 
 \frac{N \times \Frac{3}{2} k_{B} T}{V}
  \equal 
  \frac{\Frac{3}{2} p V}{V} \equal \frac{3}{2} \appear{p} \equal \appear{1.5 \times 10^{5} \mathrm{~J} / \mathrm{m}^{3}}
\]
\item So, the ratio is small: $\frac{U_{\text {photon }}}{U_{\text {matter }}}\appear{=4 \times 10^{-11}$}

 \implies{ Thus, we conclude that under NTP conditions the fraction of \\
 		\Alert[red]{thermodynamic system's total energy \appear{in the form of \\
		electromagnetic radiation} \appear{is very small}} }
\end{itemize}

\endofslide 
%\endofsection
\end{frame}
%%%%%%%%%%%%%%%

%%%%%%%%%%%%%%%
\begin{frame}
\frametitle{\texorpdfstring{\culine[\sectiontitleulinecolor]{\sectiontitle}}{\sectiontitle} }

\begin{itemize}
	\item We may interpret the integrand in the previous example for internal energy \appear{as the (total) energy of the system within a certain photon energy range} \appear{[[$E , E+d E$]]:}
\[
d U_{photon}(E, E+d E)
 \equal 
 \appear{E} \frac{1}{e^{E / k_{B} T}-1} \frac{8 \pi V}{h^{3} c^{3}} \appear{E^{2}} \appear{d E}
  \equal \frac{h f^{3}}{e^{h f / k_{B} T}-1} \frac{8 \pi V}{c^{3}} \appear{d f}
\]
\item[] this is the \CDAlert[blue]{Planck's law}, giving us the
\vspace{2.5mm}
\end{itemize}
\begin{minipage}{6.5cm}
\begin{itemize}
	\item[] spectral energy density for electromagnetic radiation or so called \CDAlert[red]{blackbody radiation}

\item The maximum peak of the curve moves linearly with temperature, following the \CDAlert[tunired]{Wien's law} 
\end{itemize}

\end{minipage}

\xdef\loc{south east}
\begin{tikzpicture}[remember picture,overlay] % anchor=south
    \node (A) [yshift=-0.25*\figYshift,xshift=\figXshift, anchor=\loc] at (current page.\loc) {\includegraphics[page=1,width=7.2cm]{Figures_booklet/SOM_9_Statistical_mechanics_figures/SOM_Figure_20.png}}; %scale=\figscale. % 8cm max hei
\end{tikzpicture}


\endofslide 
%\endofsection
\end{frame}
%%%%%%%%%%%%%%%

%%%%%%%%%%%%%%%
\begin{frame}
\frametitle{\texorpdfstring{\culine[\sectiontitleulinecolor]{\sectiontitle}}{\sectiontitle} }

\begin{itemize}
	\item Integration of the area under the curve gives the \CDAlert[red]{Stefan–Boltzmann's law}:
\[
U \propto \appear{T^{4}}
\]
\item To summarize\appear{, the \CDAlert[fuchsia]{bosonic photon gas model}} \appear{allows to derive many of the foundational equations}\appear{/laws} \appear{connecting thermodynamics} \appear{with electrodynamics}
\item[] \begin{itemize}[itemsep=2.5mm]
	 \item Planck's law
	 \item Blackbody radiation
	 \item Wien's law
	 \item Stefan–Boltzmann's law
\end{itemize}

\end{itemize}

\xdef\loc{south east}
\begin{tikzpicture}[remember picture,overlay] % anchor=south
    \node (A) [yshift=-0.25*\figYshift,xshift=\figXshift, anchor=\loc] at (current page.\loc) {\includegraphics[page=1,width=7.2cm]{Figures_booklet/SOM_9_Statistical_mechanics_figures/SOM_Figure_20.png}}; %scale=\figscale. % 8cm max hei
\end{tikzpicture}


%\endofslide 
\endofsection
\end{frame}
%%%%%%%%%%%%%%%
